\section{The Breath of Life}

\begin{quotex}
Man consists of three elements: spirit, soul, and body, which sometimes are reckoned as two, for often the soul is included in the designation of spirit (for it is that certain rational part, which beasts do not have, that is called spirit). Our chief element is the spirit. \flright{\textsc{St. Augustine}}

\end{quotex}
\textbf{Augustine} and the Eastern fathers tend to distinguish between the soul and spirit. Basing himself on Aristotle, \textbf{Aquinas} employs different terminology although the principle is the same. This can be a source of confusion; that is why we need to know our Greek and Latin.

If Adam had his spirit breathed into him by God, and Eve by implication, Fr. Warkulwiz poses the question of how his descendants got their souls. The first choice is transmission by the reproductive process. This is rejected out of hand by Augustine on the grounds that the soul would be a body rather than a spirit. However, Augustine realized that the special creation of the soul by God is inconsistent with the dogma of original sin.

Aquinas takes an important step in resolving that issue. He distinguishes between the sensitive soul which man has in common with the animals and the intellectual soul, which distinguishes man from the animals. The latter soul is the ruach that God breathed into Adam, what we have been calling “spirit”. With this distinction in mind, Aquinas can claim that that the sensitive soul is indeed transmitted through seminal generation, while the intellectual soul is not. We have already addressed the topic of heredity\footnote{\url{https://www.gornahoor.net/?p=1664}} as it applies to the body and the (sensitive) soul. (In the ensuing discussion, I will replace the idea of intellectual soul with the word “spirit”.)

But here is where the dogma gets fuzzy. He says it is unnatural for the intellectual soul (spirit) to be without a body, so he insists that the spirit is created by God and infused into the body at the moment of conception. This explanation has the benefit that it acknowledges the transcendence of the spirit, since it is not mechanically produced as part of the world process, but has a transcendent source. Nevertheless, it raises some issues that are unsatisfactory from a metaphysical and esoteric perspective.

\begin{itemize}
\item How can something immortal and transcendent have a beginning in time? 
\item If a spirit without a body is “unnatural”, how do we account for purely spiritual postmortem states? 
\item How does a perfect God create a spirit deformed by original sin? 
\end{itemize}

\paragraph{Pre-existence of spirit}
\begin{quotex}
Before I [God] formed thee in the belly I knew thee. \flright{Jeremiah 1:5}

\end{quotex}
Rene Guenon pointed out that without the ability to think outside of time, a man can never master metaphysics; hence we can't think in terms of a clock in the heavens that is synchronized to a clock on earth. To be consistent with the dogma, we could say that the spirit exists virtually, but actually when conjoined with a soul and body.

\paragraph{Postmortem states}
Here Aquinas is inconsistent. If the spirit persists in a postmortem state, then there is no reason to require it to be joined to a body, so the pre-existence of the spirit is not logically refuted. Just the opposite, since something immortal has no beginning in time. In any case, Genesis is concerned only with this world. Although Fr. Warkulwiz mentions the idea of possible worlds, he doesn't follow through with it. We have pointed out in the post on Parallel Universes\footnote{\url{https://www.gornahoor.net/?p=1863}}, that the spirit can, and does, incarnate in different worlds. This makes the question moot.

\paragraph{Incarnation}
Both \textbf{Rene Guenon} and \textbf{Valentin Tomberg} employ the symbolism of the cross. There is the horizontal line which represents the world process, or Destiny. The vertical line represents transcendence, or Will. Since God works through secondary causes, then the incarnation of a spirit should be approached in the same way. Fr. Warkulwiz points out that creation is identical with will: when we say that God creates the word, we mean he wills it into existence. Analogously, the spirit wills itself into existence at the appropriate moment in the world process. If the world process is corrupted (by original sin), then the perfect spirit will find itself in that situation.

This is something that is not proved by argument, nor by quoting texts, but rather by remembering. If I misplace my keys, it is not because I willed to lose them. Nevertheless, I am responsible for that state of affairs. The reason they are lost is because I was not acting consciously when I laid them down. As an exercise, the next time you misplace something, instead of frantically searching for it, sit quietly and try to remember where they are.

There is a similar exercise for your life. Try to remember why you chose this particular life, its qualities, your parents, your world. Then you will understand.

\flrightit{Posted on 2011-07-12 by Cologero}

\begin{center}* * *\end{center}

\begin{footnotesize}\begin{sffamily}

\texttt{Iulianus on 2011-07-13 at 07:13 said: }

Many explain the part in the Bible, “God breathed the breath of life into the face of Adam” the first-created, who was created by Him from the dust of the ground, it must mean that until that moment there was neither human soul nor spirit in Adam, but only the flesh created from the dust of the ground. This interpretation is wrong, for the Lord created Adam from the dust of the ground with the constitution which the holy Apostle Paul describes: “May your spirit and soul and body be preserved blameless at the coming of our Lord Jesus Christ” (1 Thess. 5:23). And all these parts of our nature were created from the dust of the ground, and Adam was not created dead, but an active being like all of God's animate creatures living on earth.

The point is, that if the Lord God had not breathed afterwards into his face, this breath of life – that is, the grace of our Lord God the Holy Spirit Who proceeds from the Father, rests in the Son and is sent into the world for the Son's sake – Adam would have remained without the Holy Spirit within him. It is the Holy Spirit who raised Adam to Godlike dignity. However perfect, he had been created and superior to all the other creatures of God, as the crown of creation on earth, he would have been just like all the other creatures, though they have a body, soul and spirit, each according to its kind, do not have the Holy Spirit within them. But when the Lord God breathed into Adam's face the breath of life, then, according to Moses' word, “Adam became a living soul” (Gen. 2:7), that is, completely and in every-way like God, and like Him, forever immortal. Adam was immune to the action of the elements to such a degree that water could not drown him, fire could not burn him, the earth could not swallow him in its abysses, and the air could not harm him by any kind of action whatever. Everything was subject to him as the beloved of God, as the king and lord of creation, and everything looked up to him, as the perfect crown of God's creatures. Adam was made so wise by this breath of life, which was breathed into his face from the creative lips of God, the Creator and Ruler of all, that there has never been a man on earth wiser or more intelligent, and it is unlikely that there ever will be. When the Lord commanded him to give names to all the creatures, he gave every creature a name which completely expressed all the qualities, powers and properties given it by God at its creation.

“As a result of this gift, of the supernatural grace of God, which was infused into him by the breath of life, Adam could see, understand the Lord walking in Paradise, comprehend His words, understand the conversation of the holy Angels, the language of all beasts, birds and reptiles and all that is now hidden from us the fallen and sinful creatures. All this was so clear to Adam before his fall. The Lord God also gave Eve the same wisdom, strength, unlimited power, and all the other good and holy qualities. He created her not from the dust of the ground, but from Adam's rib in the Eden of delight, the Paradise which He had planted in the midst of the earth.

“In order that they might always easily maintain the immortal, divine and perfect properties of this breath of life, God planted in the midst of the garden the tree of life with fruits endowed with all the essence and fullness of His divine breath. If they had not sinned, Adam and Eve themselves as well as all their posterity could have always eaten of the fruit of the tree of life and so would have eternally maintained the vivifying power of divine grace.

“They could have also maintained for all eternity the full powers of their body, soul and spirit in a state of immortality and perpetual youth, and they could have continued in this immortal and blessed state of theirs forever. At the present time, however, it is difficult for us even to imagine such grace.

Source: Saint Seraphim of Sarov – On Acquisition of the Holy Spirit

\url{http://www.fatheralexander.org/booklets/english/sermon\_st\_seraphim.htm}


\hfill

\texttt{Will on 2011-07-16 at 09:08 said: }

Ananda Coomaraswamy's article The Coming to Birth of the Spirit offers much clarification on the terms `soul' and `spirit,' as well as others.

The text is available online, though unfortunately without the substantial footnotes which are worthwhile in themselves.

\url{http://www.religioperennis.org/documents/acoomaraswamy/spirit.pdf}


\bigskip

\texttt{Iulianus on 2011-07-17 at 07:19 said: }

Notes on “The Coming to Birth of the Spirit“ are included in The Essential Ananda K. Coomaraswamy, the eBook can be found here: \url{http://mikrotheos.blogspot.com/2011/01/blog-post.html}

I think that St Seraphim of Sarov is right, but it should be emphasized here that “grace of our Lord God the Holy Spirit”, i.e. “the energies of God” are uncreated. Clarification from the book THEOSIS – THE TRUE PURPOSE OF HUMAN LIFE: “The energies of God are divine energies. They too are God, but without being His essence. They are God, and therefore they can deify man. If the energies of God were not divine and uncreated, they would not be God and so they would be unable to deify us, to unite us with God. There would be an unbridgeable distance between God and men. But as God has the divine energies, and unites with us by these energies, we are able to commune with Him and to unite with His Grace without becoming identical with God, as would happen if we united with His essence. So we unite with God through His uncreated energies, and not through His essence. This is the mystery of our Orthodox faith and life. Western heretics cannot accept this. Being rationalists, they do not discern between the essence; and the energy of God, so they say that they cannot speak about man's Theosis because God is only essence, for on this basis, how can man be deified when they do not accept that the divine energies are uncreated, but regard them as created? How can something created deify man, i.e. how can something outside God deify man? In order not to fall into pantheism, they do not talk about Theosis at all. What then, according to them, remains as the purpose of human life? Simply moral improvement. If man cannot be deified with divine Grace and divine energies, what purpose does his life have? Only that he becomes morally better. But moral perfection is not enough for man. It is not enough for us simply to become better than before, simply to perform moral deeds. We have as our final aim to unite with holy God Himself. This is the purpose of the creation of the universe. This is what we desire. This is our joy, our happiness, and our fulfillment.''

According to St Ignatius Brianchaninov “Our mind is – image of God the Father, our word (and unspoken word is often called “thought”) – image of the Son of God, and our spirit – image of the Holy Spirit.” 

“Whoever is joined unto the Lord is One Spirit” I Cor. 6:17

\hfill

\texttt{Cologero on 2011-07-18 at 00:55 said: }

As for the Coomaraswamy article, what specific clarification (of what has been written) are you referring to? Commaraswamy takes pains to point out that we are not discussing philosophical categories, but rather those arising from experience. Thus ultimately the only “clarification” of value comes from one's own self-understanding.

\hfill

\texttt{Cologero on 2011-07-20 at 11:53 said: }

Thanks, Iulianus, for the extensive quotes.

\hfill

\texttt{Will on 2011-07-22 at 21:17 said: }

The leading quotation from Augustine seems to need clarification, insofar as it equates the “rational” with “spirit.” Perhaps this simply reflects a problem in the translation.

Rational comes from ratio, which is discursive thinking and reasoning (dianoia in Greek). A higher faculty is Intellect – intellectus, or nous in Greek. Plato holds that man comes to know and unify with the Good through noesis, not through discursive reasoning. Does Augustine hold the contrary? I don't know; I am quite ignorant of Augustine's theology.

Coomaraswamy argues that in the trinity of Body-Soul-Spirit, it is the spirit which is eternal. Augustine says the same thing, but if he asserts that the spirit is the rational faculty, is he asserting that this faculty survives death?

My (limited) understanding of these terms is as follows: The spirit is eternal and transcendent, as well as impersonal, such that one cannot rightly speak of “my” or “your” spirit as some kind of individual thing; indeed, it is more accurately “nothing”, no-thing, as Coomaraswamy says. 

The soul, or psyche, can be spoken of, and moreover can be worked on, harmonized, ordered, and perfected, so as to fully partake and reflect the Divine. However, traditions including Christianity describe this process as soul-death, or ego-death, as when Jesus says that whoever does not hate his own soul cannot be a disciple. This seems to me to refer to the individuality of the soul vs. the non-individuality of the spirit, and the necessity to `die before you die' so as to be born again into the eternity of spirit.

Although etymologically, the use of `spirit' to denote the transcendent, eternal, and uncreated seems problematic insofar as the word comes from spirare – to breathe – a limited bodily function. Or does this choice of terms hint at the importance of the breath as a link to the transcendent?

\hfill

\texttt{Perennial on 2011-07-23 at 02:30 said: }

Will, you make an excellent point. From an English perspective, I personally prefer using the term soul for what they refer to as spirit, since in English this is generally regarded as the eternal aspect of man, and vice versa for spirit, which has a rather more nebuluous connotation. But there may be an argument for the “orthodox” use in English. Perhaps it would be that we should at least to try to use the same terminology in the same sense, or at least simili modo. In any case, I would be interested in seeing other views on this.

\end{sffamily}\end{footnotesize}
